% SPDX-License-Identifier: Apache-2.0
%
% Vreteno through an open ASIC flow. Built by //docs:tapeout. Every
% listing is included from the tree and every number comes from a run.
\documentclass[journal,9pt]{IEEEtran}

\newcommand{\listingsize}{\fontsize{6}{7}\selectfont}
% SPDX-License-Identifier: Apache-2.0
%
% The house preamble, shared by every document in this directory. Loaded
% after \documentclass, before \title. Keeping it in one file is what
% stops two articles drifting apart in font, listing style or figure
% style.
% Computer Modern. IEEEtran selects Times (ptm) by default, so the face has
% to be chosen explicitly.
%
% Latin Modern is the Computer Modern variety used here, and the choice is
% forced rather than aesthetic. Under T1 encoding, plain `cmr` has no Type 1
% outlines unless cm-super is installed, and it is not in the pinned
% distribution, so pdflatex falls back to EC bitmaps and every glyph in the
% PDF becomes a Type 3 font. That was measured: the first build produced a
% document with 10 Type 3 fonts and no Type 1. Latin Modern is Computer
% Modern redrawn with a full T1 repertoire and real .pfb outlines, so the
% same design comes out scalable.
\usepackage[T1]{fontenc}
\usepackage[utf8]{inputenc}
\usepackage{lmodern}
% microtype is not in the pinned TeX distribution. emergencystretch alone
% keeps the two-column measure from producing overfull lines.
\setlength{\emergencystretch}{3em}

\usepackage{amsmath}
\usepackage{amssymb}
\usepackage{amsthm}
\usepackage{booktabs}
\usepackage{array}
% enumitem is not in the pinned TeX distribution; the list spacing below
% does what \setlist would have done.
\usepackage{float}
\usepackage{xcolor}
\usepackage{listings}
\usepackage{tikz}
\usetikzlibrary{arrows.meta,positioning,fit,backgrounds,calc}
\usepackage[hidelinks,bookmarks=true]{hyperref}

% Numbered, definition-style box for a rule the reader is meant to apply.
\theoremstyle{definition}
\newtheorem{rulebox}{Rule}
\newtheorem{example}{Example}

% Unnumbered asides.
\theoremstyle{remark}
\newtheorem*{tip}{Tip}
\newtheorem*{pitfall}{Pitfall}

% Tighter lists, in place of enumitem's \setlist.
\setlength{\topsep}{2pt}
\setlength{\itemsep}{1pt}
\setlength{\parsep}{0pt}

\newcommand{\code}[1]{\texttt{#1}}
\newcommand{\term}[1]{\emph{#1}}

% Syntax highlighting. Four roles, distinguishable in colour and still
% distinguishable in greyscale, because an IEEE article gets printed: the
% keyword is the only bold role, the comment is the only italic one, and the
% two remaining roles differ in lightness as well as in hue.
\definecolor{kwblue}{RGB}{20,60,150}
\definecolor{cmtgreen}{RGB}{90,120,90}
\definecolor{strred}{RGB}{150,50,40}
\definecolor{numgray}{gray}{0.45}
\definecolor{framegray}{gray}{0.70}
\definecolor{bgshade}{gray}{0.97}
% A darker fill for figure cells. The listing background has to stay very
% light to keep code readable; a figure cell has to be clearly shaded or the
% distinction it draws is invisible in print.
\definecolor{cellshade}{gray}{0.90}

% The listing size is the document's choice, because it depends on the
% column: a two-column article sets `\listingsize` to 6pt and keeps its
% listings within 70 columns, a one-column document keeps the body size
% and includes source files at 80. Lines are never broken; a line that does not fit
% overflows the frame, which is visible, rather than wrapping, which is
% not.
\providecommand{\listingsize}{\scriptsize}

% `'` is deliberately not a string delimiter: the language uses it for loop
% labels, as Rust does, so treating it as a quote would swallow the rest of
% the line.
\lstdefinestyle{txhdl}{
  basicstyle=\ttfamily\listingsize,
  keywordstyle=\bfseries\color{kwblue},
  commentstyle=\itshape\color{cmtgreen},
  stringstyle=\color{strred},
  numberstyle=\tiny\color{numgray},
  morekeywords={mod,use,pub,const,type,struct,enum,trait,impl,fn,async,let,mut,
    match,if,else,loop,while,for,in,break,continue,return,as,await,
    module,bus,channel,transaction,protocol,pipeline,flow,seq,config,
    port,stage,pipe,instance,connect,bind,tag,domain,blackbox,foreign,
    property,role,signal,handler,true,false,
    sequence,interface,modport,implements,package,import,var,wire,func,proc,
    case,default,next,fork,branch,join,state,fsm},
  morecomment=[l]{//},
  morestring=[b]",
  showstringspaces=false,
  columns=fullflexible,
  keepspaces=true,
  breaklines=false,
  % Framed, so a listing is bounded rather than floating in the column.
  frame=single,
  rulecolor=\color{framegray},
  framesep=3pt,
  backgroundcolor=\color{bgshade},
  xleftmargin=3pt,
  xrightmargin=3pt,
  aboveskip=6pt,
  belowskip=6pt,
  % Every listing is numbered and captioned, so it can be referred to by
  % number. `caption` without `float` numbers the listing and leaves it
  % where it was written, which is what a listing in running prose needs.
  captionpos=b,
  abovecaptionskip=2pt,
  belowcaptionskip=6pt,
}

% Figures drawn as text. Framed the same way, and with the highlighting
% switched off, because a box-and-arrow diagram is not code and half its
% words turning blue reads as an accident. Setting a style adds keys rather
% than clearing the ones already set, so the three roles are neutralised by
% hand rather than by leaving them out.
\lstdefinestyle{ascii}{
  basicstyle=\ttfamily\listingsize,
  keywordstyle=\ttfamily,
  commentstyle=\ttfamily,
  stringstyle=\ttfamily,
  columns=fullflexible,
  keepspaces=true,
  frame=single,
  rulecolor=\color{framegray},
  framesep=3pt,
  backgroundcolor=\color{bgshade},
  xleftmargin=3pt,
  xrightmargin=3pt,
  aboveskip=6pt,
  belowskip=6pt,
}
\lstset{style=txhdl}

% House TikZ styles. `shorten` lives on the arrow styles rather than on each
% arrow, so every arrow in the document gets the gap, including ones added
% later. TikZ clips a path at a node's border, so without it an arrowhead
% butts against the box with no gap at all. Keep `node distance` well above
% twice the shorten, or the arrow shrinks to nothing.
\tikzset{
  box/.style={draw,rounded corners=2pt,align=center,inner sep=4pt,
              font=\footnotesize},
  boxf/.style={box,fill=bgshade},
  lbl/.style={font=\scriptsize\itshape,align=center},
  fl/.style={-{Stealth[length=4pt]},shorten >=2.5pt,shorten <=2.5pt},
  fld/.style={fl,dashed},
}



\InputIfFileExists{buildstamp.tex}{}{}
\providecommand{\buildstamp}{Draft build (outside the Bazel flow).}

% The numbers the prose quotes, written by the run that produced them.
\input{asic_counts.tex}

% A script of the flow, highlighted for its own vocabulary: the
% comment character is the hash, and the commands are the steps.
\lstdefinestyle{script}{%
  style=txhdl,
  morecomment=[l]{\#},
  morekeywords={set,expr,proc,foreach,source,puts,format,open,close,
    read_verilog,hierarchy,synth,opt,dfflibmap,abc,setundef,splitnets,
    opt_clean,hilomap,insbuf,tee,stat,write_verilog,read_lef,
    read_liberty,read_sdc,link_design,initialize_floorplan,place_pins,
    tapcell,pdngen,global_placement,detailed_placement,check_placement,
    clock_tree_synthesis,set_propagated_clock,repair_design,
    repair_clock_nets,repair_timing,global_route,detailed_route,
    filler_placement,estimate_parasitics,write_def,create_clock,
    set_input_delay,set_output_delay,set_driving_cell,set_load,
    set_max_fanout,set_clock_uncertainty,set_thread_count,get_ports,
    get_clocks,all_inputs,all_outputs,current_design,lsearch,metric,
    report_design_area,report_worst_slack,report_tns,report_checks,
    report_power,report_check_types,report_clock_skew},
}

\title{A Tapeout Prep in TxHDL: Vreteno on an Open 45\,nm Node}

\author{Filip Filmar and Dragiša Janković%
\thanks{Both authors are independent. Filip Filmar is the
corresponding author, at f@hdlfactory.com; Dragiša Janković is
reached at linkedin.com/in/dragisa-jankovic.}%
\thanks{This is the tutorial on what a design does after synthesis
and before a mask. The core is the one in~\cite{vreteno}; the
netlist it starts from is the one the lowering emits. Every listing
is included from the project repository \code{HDL/txhdl} at build
time. The cover~\cite{cover} lists every document.}%
\thanks{The authors used a large language model, Claude, as an
assistant in exploring the concepts and in writing and constructing
the documents and the programs; every commit of the repository says
so and carries the prompts that led to it, verbatim.}%
\thanks{\buildstamp}}

\markboth{A Tapeout Prep in TxHDL}{Filmar and Janković: A Tapeout Prep in TxHDL}

\begin{document}
\maketitle

\begin{abstract}
A design that has been simulated, lowered, checked against its trace
and synthesised onto an FPGA has still never been a shape. This takes
the Vreteno core's Verilog, the same file the board's bitstream is
built from, maps it onto an open 45\,nm standard cell library with
Yosys, and floorplans, places, routes and times it with OpenROAD. The
result is a die \asicdie\,$\mu$m square holding \asiccells{} cells
that meets a \asicperiod\,ns clock with \asicslack\,ns to spare and
routes with \asicdrc{} violations. Both tools are fetched by
checksum and unpacked by the build, so the flow needs nothing
installed. This is not a tapeout, and the last section says what a
tapeout would still want.
\end{abstract}

\section{What This Is}
\label{sec:t-what}

An FPGA answers one question about a design: does it fit, and how
fast does it go, on this part. An ASIC flow answers a different one,
and asks the design for things an FPGA never does. There is no block
RAM to put a memory in and no arithmetic block to put a multiplier
in, so both become gates. There is no configuration bitstream to give
a register its value at power on, so a register that was initialised
and never reset is a register with no defined value. And there is no
routing fabric that exists whether or not the design uses it: every
wire is drawn, and its length is a delay.

The Vreteno core has been through the first of those~\cite{vreteno}.
It runs on an Alinx AX7A200 at 100\,MHz, and what it said on the
board is what its simulation said. This document puts the same
netlist through the second, and reports what the flow found.

None of it is a tapeout. No foundry has a process that matches this
library, no mask was made, and the checks a foundry demands are not
here. Section~\ref{sec:t-not} lists what is missing, one item at a
time, because a document that showed a picture of a layout and
stopped would be claiming more than it had done.

\section{The Node}
\label{sec:t-node}

Nangate45 is an open standard cell library at a notional 45\,nm: a
timing model for each cell, a layout abstraction saying what shape it
is and where its pins are, and a technology file saying what the ten
metal layers are and how wide a wire on each may be. It is the
library the OpenROAD project uses for its own examples, it is
published under Apache 2.0 with the flow scripts, and the build
fetches the six files the flow reads, each by checksum, at one
revision, with its licence.

It is not a foundry's library. No wafer was ever made from it, its
numbers are a plausible 45\,nm and not a measured one, and it has no
design rule deck that a fab would sign. That is exactly why it is
the right thing to put a first flow on: every step is real, every
number means what it says, and nothing is under a licence that keeps
it out of a repository.

The one thing it does not have is memory. A foundry node comes with
a compiler that makes an SRAM macro of whatever shape is asked for,
and there is no such thing here; the flow scripts ship mock macros of
fixed shapes for experimenting with placement, and a mock macro has
no contents. Section~\ref{sec:t-cells} says what the core's memories
became instead, which is the most interesting thing this flow found.

\section{The Tools, and How They Are Pinned}
\label{sec:t-tools}

Everything in this repository is fetched by checksum and nothing is
installed, and the ASIC flow is held to the same rule. That is
harder here than elsewhere, because Yosys and OpenROAD are published
as Debian packages, and a Debian package names the shared libraries
it needs without carrying them.

So the build resolves the closure first. \code{//tools/debclosure}
reads a suite's package index, walks the dependencies from the
packages named on its command line, and writes a lock file: one row
per package with its version, its size, its checksum and two places
to fetch it from. The first is the live archive, which is fast and
forgets a package as soon as a newer one replaces it; the second is
\code{snapshot.debian.org}, which is slow and never forgets. A
fetcher that tries them in order is fast while the pin is current and
still right years later.

Two such lock files are in the tree. Yosys 0.52 and the ten libraries
it needs are eleven packages and thirteen megabytes. OpenROAD is the
last release its project published as a package, fifty two megabytes
of its own, and the closure of what it links against is a hundred and
thirty one packages and seventy two megabytes more, most of that Qt,
which the program links against whether or not a window is ever
opened. Beside them is the solver its placer uses, which the package
names in its runtime path and does not carry.

Nothing unpacks them but Bazel. A \code{.deb} is an \code{ar} archive
holding a \code{tar.xz}, and Bazel's own decompressors open both, so
the repository rule in \code{//third\_party/debs} downloads each
package, unpacks it twice, and leaves one tree that the two tools run
against with their own libraries. There is no \code{dpkg}, no
container, and no step that reads the machine.

\section{From Verilog to Cells}
\label{sec:t-cells}

Listing~\ref{lst:t-synth} is the whole of the mapping. The build
fills in the words between the at signs, and the rest is the order
the OpenROAD flow scripts use.

\lstinputlisting[style=script, caption={Vreteno onto standard
cells.}, label={lst:t-synth}]{cpu/vreteno/asic/synth.ys}

Six of those lines are not decoration. \code{dfflibmap} must run
before \code{abc}, because a register can only be mapped to a cell
that exists and the combinational mapper will not invent one.
\code{setundef} gives a defined value to bits that had none, and
\code{splitnets} turns every bus into its bits, because a netlist of
cells has no such thing as a vector. \code{hilomap} replaces the
constants a netlist writes as one and zero with the two cells that
drive them, which the mapper will not use on its own because the
library marks them \code{dont\_use}. And \code{insbuf} puts a buffer
on a port wired straight through to another, because a layout has no
such thing as a wire with no driver.

Their order matters more than it looks. Buffering before the buses
are split puts a buffer on every bit of every bus: three thousand
eight hundred cells that nothing needed, a fifth of the design, found
by reading the cell counts and not by any tool complaining.

What comes out is \asiccells{} cells on \asicnets{} nets, of which
\asicflops{} are flip-flops.

\subsection{The Memories}

Compare that with the same core on the Artix-7, which is 2151 lookup
tables, 489 flip-flops, two block RAM tiles and four arithmetic
blocks~\cite{vreteno}. The flip-flop counts differ by a factor of
three, and the reason is the register file. On the FPGA it is twelve
distributed memory primitives, because the fabric has a small
memory that is not a register; on an ASIC there is no such thing, so
thirty two words of thirty two bits are a thousand and twenty four
flip-flops. Those and the core's own registers come to \asicflops{},
against 489 on the board, where the register file was not flip-flops
at all.

The instruction memory went the other way, and this is the
interesting part. On the FPGA it is two block RAM tiles of 1024 words
by 32 bits, one for each of the two words the fetch reads, since an
instruction may cross from one word into the next. Here there is no
block RAM, and the memory is read-only:
nothing in the core writes it, and its contents are the program,
which is known at synthesis. So it is not a memory at all. It is a
mask ROM, which is to say a lump of combinational logic, and what it
costs depends on the program in it. A different program is a
different area. That is true of any mask ROM and it is worth saying
out loud, because it means the area figure below is the area of this
core running this program, and nothing more general.

The multiplier is the third case. Four arithmetic blocks on the
FPGA, and here the gates that a thirty two bit multiply step is,
because the library has cells and nothing else.

\section{From Cells to a Layout}
\label{sec:t-pnr}

Listing~\ref{lst:t-sdc} is what the design is asked to meet. Three
nanoseconds is 333\,MHz, three times what the board runs at, because
a 45\,nm standard cell is a faster thing than an FPGA lookup table
and the point of naming a period is that every number after it means
something. The ports are timed as if a chip of the same speed were on
the other side of each, with a fifth of the period to arrive and a
fifth to be taken. Ports with no timing at all are the usual way a
report comes out clean and wrong.

\lstinputlisting[style=script, caption={What the core is asked to
meet.}, label={lst:t-sdc}]{cpu/vreteno/asic/vreteno.sdc}

The flow itself is in \code{cpu/vreteno/asic/pnr.tcl} and runs in
this order. A floorplan: a square of silicon with room for the cells,
the rows they will stand in, and the tracks they will be wired along.
The ports become pins on the edge, on the two upper layers that are
clear of the cells. Cells that tie the wells to the supplies go in,
one every 120 microns along a row and one at each end. A power grid
goes over the whole of it: the rails the rows sit on, and straps
above them on two upper layers. Then the cells are spread out,
resized where a signal was too slow or too sharp, and packed into the
rows. Then the clock: a tree of buffers so that it reaches every one
of the flip-flops at nearly the same moment, after which the
registers move a little to make that tree shorter. Then routing,
first as a plan and then as metal, and a last pass of resizing
against the delays the real wires turned out to have. Last of all,
filler cells, which carry the wells and the supply rails across the
gaps left in the rows and which come last because they have no signal
pins and so nothing to route.

\begin{figure*}[t]
\centering
\input{asicmap.tex}
\caption{The placed core, as a grid of \asicbins{} squares a side over
the die. On the left, how much of each square is filled by cells of
any kind, before the filler goes in: the pale border is the margin
the floorplan leaves outside the rows. On the right, the flip-flops
alone, shaded against the fullest square, which is \asicffpeak\,\% of
its area; they are a tenth of the cells, so against a full square the
whole panel would be one pale wash. The state of the design has a
shape: where there are no flip-flops at all there is combinational
logic, and by far the largest piece of that is the instruction
memory, which is a mask ROM and holds no state.}
\label{fig:t-map}
\end{figure*}

\section{What the Run Says}
\label{sec:t-results}

Table~\ref{tab:t-results} is the run. Every number in it was written
by the flow into a file, and the prose here quotes that file through
macros, so neither can drift from the other.

\begin{table}[htbp]
\caption{Vreteno on Nangate45, after routing.}
\label{tab:t-results}
\centering
\begin{tabular}{@{}lr@{}}
\toprule
Cells, after mapping & \asiccells{} \\
Nets & \asicnets{} \\
Flip-flops & \asicflops{} \\
Instances placed, before filler & \asicinstances{} \\
\addlinespace
Die & \asicdie{} $\times$ \asicdie\,$\mu$m \\
Cell area & \asicarea\,$\mu$m$^2$ \\
Core filled by cells & \asicutil\,\% \\
\addlinespace
Clock & \asicperiod\,ns (\asicfreq\,MHz) \\
Worst setup slack & \asicslack\,ns \\
Worst hold slack & \asicholdslack\,ns \\
Total negative slack & \asictns\,ns \\
\addlinespace
Routing violations & \asicdrc{} \\
\bottomrule
\end{tabular}
\end{table}

The timing closes, and it closes by a margin thin enough to be worth
looking at: \asicslack\,ns out of \asicperiod. Total
negative slack is \asictns, so the worst path is the only path that
is close. Hold is met everywhere as well, which matters more than it
sounds: a setup failure is a chip that has to be run slower, and a
hold failure is a chip that does not work at any speed.

Routing finished with \asicdrc{} violations. Zero is the only
interesting value that number can have, because a routed design with
violations is a design that has not been routed.

The utilization deserves a word, because \asicutil\,\% is loose. The
floorplan is sized from the cell area the mapping came to, and the
resizer then grows the design by a quarter; a production flow
iterates on that, and this one does not. A tighter
die would be smaller and harder to route, and the point of this run
is the flow rather than the squeeze.

\section{Where This Sits in the Build}
\label{sec:t-build}

Two targets. \code{//cpu/vreteno/asic:vreteno\_cells} is the mapping.
It takes about ten seconds and is built with everything else, so a
change to the core that the library cannot express is a broken build
on the next commit rather than a surprise months later.
\code{//cpu/vreteno/asic:vreteno\_pnr} is the layout. It takes about
twenty minutes, so it is tagged \code{manual}, as every Vivado target
is, and \code{bazel test //...} stays a thing somebody runs before
every commit.

That leaves every number in this document coming from a run the build
does not make, which is the problem the project's timeline has, and
it is handled the same way. The run writes two small files, the
numbers and a grid of where the cells went, and those two are
committed; \code{//tools/denmap} turns them into the table's macros
and into Fig.~\ref{fig:t-map} at build time. The routed design itself
is forty megabytes of DEF and stays where it was made. So the figure
is never a picture somebody pasted in, the prose can only quote
numbers that a tool wrote, and moving any of them means running the
layout again and committing what it says.

\section{What a Tapeout Would Still Want}
\label{sec:t-not}

This flow ends with a routed design that meets its timing. A tapeout
begins there.

\emph{A foundry, and its rules.} Nangate45 has no design rule deck
and no device models, so the routing checks are the router's own
rules and not a fab's. Physical verification, which is design rule
checking and layout versus schematic against the foundry's decks, has
not been done because there is nothing to do it against. On an open
node that has a fab behind it, Sky130 or GF180, both exist and both
would run.

\emph{Something to hold the program.} The instruction memory is a
mask ROM here, which means the program is fixed when the mask is. A
chip that is to be useful wants the program loaded, which means an
SRAM macro from the node's compiler, or a serial port that writes
into one before the core is let go. The core already has the bus for
it.

\emph{Pads.} The design has \asiccells{} cells and no pad ring. Every
port here is a pin on the edge of the core, which is what the next
level up would connect to; a chip needs pads with their drivers,
their electrostatic protection and their bond or bump pattern, and a
power ring to feed them.

\emph{A way to test it.} There is no scan chain. Every flip-flop in
the design is an ordinary one, so a manufactured part could be run
but not shifted through, and a fault in the middle of it could not be
found. Inserting scan is a synthesis option and a set of extra pins,
and it is the difference between a chip that works and a chip you can
show works.

\emph{Reset.} The core resets its pipeline, and the register file is
not reset at all. On an FPGA that is harmless: the bitstream writes
every register on the way in. On silicon the register file powers up
holding nothing in particular, which the instruction set does not
mind for thirty one of its thirty two words and does mind for the
one that must read as zero. This core reads that one from a constant
rather than from the file, so it survives, and the general point
stands anyway: an initial value in Verilog is a simulation's
convenience and not a circuit.

\emph{Corners.} One timing library was read: typical process, one
voltage, one temperature. A signoff wants the slow and the fast
corner as well, and hold checked at the fast one, because a chip has
to work at both ends of what a wafer does.

\emph{Extraction.} The delays here are estimated from the routing.
Real signoff extracts the resistance and capacitance of the metal
that was actually drawn and times against that.

Every one of those is a known piece of work rather than an open
question, and the flow they would attach to is the one above.

\begin{thebibliography}{9}

\bibitem{vreteno}
F.~Filmar and D.~Janković, ``Vreteno: an RV32IMC core in TxHDL,''
project repository \code{HDL/txhdl}, \code{//docs:vreteno}, 2026.

\bibitem{cover}
F.~Filmar and D.~Janković, ``TxHDL documents,'' project repository
\code{HDL/txhdl}, \code{//docs:cover}, 2026.

\end{thebibliography}

\end{document}
